\chapter{Introduction}

AstraZeneca PLC is a global, science-led pharmaceutical company, listed in London, Stockholm and New York \cite[p.~71]{AR}, with \$58,739m of Total Revenue \cite[p.~125]{AR} and 96,100 employees in 2025 \cite[p.~27]{AR}. This section presents what the company is, what it does, where it operates and how large it is, and closes with its contribution to the UN Sustainable Development Goals.

{\small\color{azmuted}
Unless stated otherwise, the data come from AstraZeneca's \textit{Annual Report and Form 20-F Information 2025} (``the Annual Report'', abbreviated ``AR'' in the text) \cite{AR}; page numbers are the ones printed on the report's pages. AstraZeneca PLC and its subsidiaries are called ``the Group'', values are in US dollars, and ``\$m'' and ``\$bn'' mean millions and billions. Blocks marked Analysis are interpretation by the authors, not statements from the Annual Report.
\par}

\section{Type of organisation}

AstraZeneca PLC is a public limited company, limited by shares, incorporated and domiciled in England and Wales \cite[p.~199]{AR}. It was incorporated on 17 June 1992 under the Companies Act 1985, and its registered office and corporate headquarters are in Cambridge, United Kingdom \cite[Item~4.A]{F20}, \cite[p.~230]{AR}. In Portuguese terms it corresponds to a listed \textit{sociedade anónima}, a legal form the course material cites among private companies \cite[p.~5]{INT}.

\begin{itemize}[leftmargin=*]
\item \textbf{1992--1993.} Company incorporated in 1992. In 1993 the pharmaceutical, agrochemical and specialty chemical businesses of Imperial Chemical Industries PLC were demerged; the company was called Zeneca Group PLC until April 1999. \cite[Item~4.A]{F20}
\item \textbf{1999.} Merger with Astra of Sweden; the company changed its name to AstraZeneca PLC. \cite[Item~4.A]{F20}
\item \textbf{2000.} The agrochemical business was demerged and merged with a similar business of Novartis to form Syngenta AG. \cite[Item~4.A]{F20}
\item \textbf{2007.} Acquisition of MedImmune, a biologics and vaccines business based in the United States. \cite[Item~4.A]{F20}
\item \textbf{2021.} Acquisition of Alexion, a rare disease business based in the United States. \cite[Item~4.A]{F20}
\end{itemize}

The course material classifies organisations by who benefits from them. AstraZeneca falls under the \textbf{commercial-interest} type, which is usually a private firm whose main benefit goes to its owners \cite[pp.~4--5]{INT}, and under the \textbf{private} type when classified by who holds the capital \cite[p.~6]{INT}. It is a \textbf{for-profit} company: in 2025 it made a profit for the period of \$10,233m and paid dividends of \$4,846m \cite[p.~125]{AR}; for 2025 the Board declared a full-year dividend of \$3.20 per share \cite[p.~3]{AR}, estimated at \$4,963m in total \cite[p.~225]{AR}.

Its shares are listed on the London Stock Exchange, Nasdaq Stockholm and the New York Stock Exchange \cite[p.~71]{AR}. Since 2 February 2026 the ordinary shares are listed directly on the NYSE, replacing the earlier US listing of American Depositary Shares on Nasdaq \cite[p.~50]{AR}.

\section{Capital holders}

On 31 December 2025 the share capital comprised 1,550,907,927 ordinary shares of \$0.25 each, plus 50,000 redeemable preference shares \cite[p.~224]{AR}. The AR states that major shareholders do not have different voting rights \cite[p.~225]{AR} and describes a geographically diverse shareholder base \cite[p.~76]{AR}. Four groups held 3\% or more.

\begin{xltabular}{\linewidth}{@{}Yrrr@{}}
\caption{Shareholders with 3\% or more of the voting rights. Wellington Management Company LLP, a subsidiary, is included in the Wellington Management Group figure. Source: \cite[p.~225]{AR}.}\label{tab:acionistas}\\*
\toprule
\textbf{Shareholder} & \textbf{Shares at last disclosure} & \textbf{\% at last disclosure} & \textbf{\% at 31 Dec 2025} \\*
\midrule
\endfirsthead
\toprule
\textbf{Shareholder} & \textbf{Shares at last disclosure} & \textbf{\% at last disclosure} & \textbf{\% at 31 Dec 2025} \\*
\midrule
\endhead
\bottomrule
\endfoot
BlackRock & 100,885,181 & 6.96 & \textbf{6.50} \\
The Capital Group Companies & 77,125,348 & 4.97 & \textbf{4.97} \\
Wellington Management Group LLP & 65,120,892 & 4.96 & \textbf{4.20} \\
Investor AB & 51,587,810 & 3.93 & \textbf{3.33} \\
\end{xltabular}

The largest holding at 31 December 2025 is therefore 6.50\%, and the Company has not been notified of any other person holding a notifiable interest.

\section{Activities and scope}

AstraZeneca describes itself as a global, science-led and patient-focused pharmaceutical business, committed to excellence in the research, development, manufacturing and commercialisation of prescription medicines across primary care, specialty care and rare diseases \cite[p.~8]{AR}.

The commercial organisation is split into three regions: the US, Europe-Canada and International, which includes Australia, New Zealand and the Emerging Markets (China among them). Japan reports separately \cite[p.~32]{AR}. Products reach the market through wholly-owned local marketing companies, distributors and local representative offices \cite[p.~32]{AR}.

\begin{itemize}[leftmargin=*]
\item \textbf{Research and development --- 16,100:} employees in R\&D across the Group's global sites \cite[p.~27]{AR}
\item \textbf{Commercial --- 47,400:} commercial employees; the network has 10 global hubs \cite[p.~27]{AR}
\item \textbf{Manufacturing --- 16,900:} employees in manufacturing, in 31 operations sites in 15 countries \cite[p.~27]{AR}
\item \textbf{Total --- 96,100:} employees (2024: 94,300) \cite[p.~27]{AR}
\end{itemize}

\section{Main products}

Revenue is organised in three therapy areas. Oncology is the largest, followed by BioPharmaceuticals, which groups cardiovascular, renal and metabolism, respiratory and immunology, and vaccines and immune therapies \cite[pp.~2, 4]{AR}. Rare Disease comes third.

\begin{table}[H]\centering
\caption{Total Revenue by therapy area in 2025. Percentages add to 101\% because of rounding. Source: \cite[p.~2]{AR}.}
\label{tab:receita-area}
\begin{tabular}{@{}lrrB@{}}
\toprule
\textbf{Therapy area} & \textbf{\$bn} & \textbf{\%} & \\
\midrule
Oncology            & 25.6 & 44 & \barra{1.000} \\
BioPharmaceuticals  & 23.0 & 39 & \barra{0.898} \\
Rare Disease        &  9.1 & 16 & \barra{0.355} \\
Other Medicines     &  1.0 &  2 & \barra{0.039} \\
\midrule
\textbf{Total}      & \textbf{58.7} & \textbf{100} & \\
\bottomrule
\end{tabular}
\end{table}
The AR gives the sales of the best-selling brands. Together, the five largest products brought in \$30,045m, which is 51.1\% of Total Revenue (calculated from \cite[p.~56]{AR}: 30,045 ÷ 58,739). The Group counts 16 products above \$1bn in annual sales \cite[p.~50]{AR}.

\begin{table}[H]\centering
\caption{Five largest-selling products in 2025. Source: \cite[p.~56]{AR}.}
\label{tab:top5}
\begin{tabular}{@{}lrB@{}}
\toprule
\textbf{Product} & \textbf{\$m} & \\
\midrule
Farxiga   & 8,492 & \barra{1.000} \\
Tagrisso  & 7,254 & \barra{0.854} \\
Imfinzi   & 6,063 & \barra{0.714} \\
Ultomiris & 4,718 & \barra{0.556} \\
Calquence & 3,518 & \barra{0.414} \\
\bottomrule
\end{tabular}
\end{table}

\begin{table}[H]\centering
\caption{Best-selling brands by therapy area. The AR links the fall in Soliris to patients switching to Ultomiris \cite[pp.~24, 56]{AR}.}
\label{tab:marcas}
\begin{tabular}{@{}llrr@{}}
\toprule
\textbf{Area} & \textbf{Product} & \textbf{\$m} & \textbf{Change} \\
\midrule
Oncology \cite[p.~13]{AR}     & Tagrisso  & 7,254 & +10\% \\
                              & Imfinzi   & 6,063 & +29\% \\
                              & Calquence & 3,518 & +12\% \\
                              & Lynparza  & 3,279 & $-$11\% \\
\addlinespace
Rare Disease \cite[p.~23]{AR} & Ultomiris & 4,718 & +20\% \\
                              & Soliris   & 1,837 & $-$29\% \\
                              & Strensiq  & 1,678 & +19\% \\
                              & Koselugo  &   662 &  +5\% \\
\bottomrule
\end{tabular}
\end{table}
\section{Main markets}

The United States and Europe together account for 65\% of Total Revenue \cite[p.~2]{AR}. The Emerging Markets are the second-largest region and the fastest growing. ``Established RoW'' covers Japan, Canada, Australia and New Zealand \cite[p.~33]{AR}.

\begin{table}[H]\centering
\caption{Total Revenue by region in 2025, with growth at actual exchange rates. The shares are 43\%, 26\%, 22\% and 9\% \cite[p.~2]{AR}; the values come from \cite[p.~32]{AR} and add to \$58,739m.}
\label{tab:receita-regiao}
\begin{tabular}{@{}lrrB@{}}
\toprule
\textbf{Region} & \textbf{\$m} & \textbf{Growth} & \\
\midrule
United States    & 25,450 & +10\% & \barra{1.000} \\
Emerging Markets & 15,303 & +12\% & \barrab{0.601} \\
Europe           & 12,739 &  +5\% & \barra{0.501} \\
Established RoW  &  5,247 &  +5\% & \barrab{0.206} \\
\bottomrule
\end{tabular}
\end{table}
\section{Countries of presence}

AstraZeneca has an active presence in more than 80 countries and sells its products in more than 125 \cite[p.~32]{AR}. Its network includes six strategic R\&D centres (Boston, Gaithersburg, Cambridge, Gothenburg, Shanghai and Beijing, the latter opened in 2025), ten global hubs, and 31 operations sites in 15 countries \cite[p.~27]{AR}.

\figura{fig04-alcance-global}{1.0}{Global reach and presence: R\&D centres, global hubs, operations sites and employees by reporting region. Source: \cite[p.~27]{AR}.}

\begin{table}[H]\centering
\caption{Employees by reporting region at the end of 2025, out of a total of 96,100. Source: \cite[p.~27]{AR}.}
\label{tab:empregados}
\begin{tabular}{@{}lrrB@{}}
\toprule
\textbf{Region} & \textbf{Employees} & \textbf{\%} & \\
\midrule
Emerging Markets & 36,800 & 38 & \barra{1.000} \\
Europe           & 34,000 & 35 & \barrab{0.924} \\
United States    & 17,900 & 19 & \barra{0.486} \\
Established RoW  &  7,400 &  8 & \barrab{0.201} \\
\bottomrule
\end{tabular}
\end{table}
\section{Key indicators}

\begin{table}[H]\centering
\caption{Key indicators for 2025. Sources: \cite[p.~1]{AR}; \cite[p.~125]{AR}; \cite[p.~27]{AR}.}
\label{tab:indicadores}
\begin{tabular}{@{}lr@{}}
\toprule
\textbf{Indicator (2025)} & \textbf{Value} \\
\midrule
Total Revenue                   & \$58,739m (+9\%) \\
\quad Product Sales             & \$55,573m \\
\quad Alliance Revenue          & \$3,067m \\
\quad Collaboration Revenue     & \$99m \\
\addlinespace
R\&D expense                    & \$14,232m \\
Operating profit (reported)     & \$13,743m \\
Profit for the period           & \$10,233m \\
Dividends declared and paid     & \$4,846m \\
Employees                       & 96,100 \\
\bottomrule
\end{tabular}
\end{table}

\begin{table}[H]\centering
\caption{Total Revenue from 2023 to 2025, in \$bn. Source: \cite[p.~11]{AR}.}
\label{tab:receita-3anos}
\begin{tabular}{@{}lrB@{}}
\toprule
\textbf{Year} & \textbf{\$bn} & \\
\midrule
2023 & 45.8 & \barra{0.780} \\
2024 & 54.1 & \barra{0.920} \\
2025 & 58.7 & \barra{1.000} \\
\bottomrule
\end{tabular}
\end{table}
\section{Contribution to the UN Sustainable Development Goals}

The AR mentions the UN SDGs once, when it says that partnerships with multilateral organisations and NGOs support them \cite[p.~76]{AR}. It does not map actions to individual goals. The mapping below is therefore Analysis by the authors: each goal is linked to the actions and targets the Group reports.

\subsection{What the Group already does}

\begin{xltabular}{\linewidth}{@{}c>{\hsize=1.45\hsize}Y>{\hsize=0.55\hsize}Y@{}}
\caption{SDGs already supported by the Group.}\label{tab:ods-atuais}\\*
\toprule
\textbf{SDG} & \textbf{Action and evidence} & \textbf{Why it contributes} \\*
\midrule
\endfirsthead
\toprule
\textbf{SDG} & \textbf{Action and evidence} & \textbf{Why it contributes} \\*
\midrule
\endhead
\bottomrule
\endfoot
3 & \textbf{Good health and well-being.} 320 million people positively impacted since 2024, 156 million of them from underserved groups. Ambition for 2030: one billion people, 400 million from underserved communities. A global tiered pricing policy with four tiers based on gross national income, and Patient Assistance Programmes. \cite[pp.~5, 41]{AR} & The core business is to discover and supply medicines, and these programmes widen access to treatment. \\
13 & \textbf{Climate action.} Scope 1 and 2 emissions are down 88.1\% since 2015, against a target of 98\% by 2026. Targets: $-$50\% of Scope 3 by 2030 and $-$90\% by 2045 (2019 baseline), and net zero by 2045. \cite[pp.~2, 42]{AR} & It sets measurable reduction targets, verified under the Science Based Targets initiative \cite[p.~42]{AR}. \\
6 & \textbf{Clean water and sanitation.} The Group's website reports a 23\% reduction in water withdrawal against 2015 \cite{NATURE}. The AR gives a 24\% reduction in water use for the Operations function alone \cite[p.~33]{AR}. & It reduces the use of a shared resource in its own operations. \\
12 & \textbf{Responsible consumption and production.} Waste is down 23\% since 2015, with a goal of a 50\% waste circularity rate by 2030 \cite{NATURE}. The AR also gives $-$23\% waste for Operations \cite[p.~33]{AR}. & It cuts waste at source and sets a circularity goal. \\
15 & \textbf{Life on land.} AZ Forest, a reforestation and biodiversity programme: 66 million trees planted on six continents since 2020 (at the end of 2025). The website states a target of more than 100 million by 2030 and that the target was revised in 2026 \cite{FOREST}. & It restores ecosystems and biodiversity. \\
\end{xltabular}

\subsection{Additional goals supported by AR evidence}

\begin{xltabular}{\linewidth}{@{}cY@{}}
\caption{Further SDGs supported by evidence in the Annual Report.}\label{tab:ods-adicionais}\\*
\toprule
\textbf{SDG} & \textbf{Evidence} \\*
\midrule
\endfirsthead
\toprule
\textbf{SDG} & \textbf{Evidence} \\*
\midrule
\endhead
\bottomrule
\endfoot
5 & \textbf{Gender equality.} Women are 54\% of the workforce, there are seven women on the Board (50\%), and five of the ten members of the Senior Executive Team are women \cite[p.~39]{AR}. \\
17 & \textbf{Partnerships for the goals.} More than 1,000 collaborations worldwide \cite[p.~37]{AR}, and partnerships with multilateral organisations and NGOs in support of the SDGs \cite[p.~76]{AR}. \\
\end{xltabular}

\subsection{What the Group could do}

The AR itself shows where the biggest gaps are. The three proposals below start from those gaps and are not actions the Group has announced.

\begin{analysis}
\begin{itemize}[leftmargin=*,nosep]
\item \textbf{SDG 13: focus on suppliers.} More than 95\% of the Group's greenhouse-gas emissions are in the upstream and downstream value chain, and it already engages suppliers on science-based targets \cite[p.~42]{AR}. Strengthening supplier decarbonisation would act where most emissions are.
\item \textbf{SDG 3: close the equity gap.} The ambition is to reach 400 million people from underserved communities by 2030, and 156 million had been reached since 2024 \cite[pp.~5, 41]{AR}. Concentrating access programmes on underserved groups would accelerate progress towards that target.
\item \textbf{SDG 6, 12 and 15: environmental impact of medicines.} The AR states that pharmaceutical residues entering the environment are currently an unavoidable result of patient use of medicines, and reports measures such as environmental risk assessments \cite[p.~45]{AR}. Extending those measures along the supply chain is a possible next step.
\end{itemize}
\end{analysis}
